\begin{lemma}[Extension-base Part I(2): delta-color exclusion]
Use the setup of \noderef{ExtensionBaseSupportOrder}: the terminal
iteration \(n\) is a PE with support \(v_n\), \(j=m(v_n)\), and \(L\)
is the active set.  Let the Lean index called \(jSlot\) represent the
manuscript stage \(j\), so its coloring and tree are
\(\varphi_{j-1}\) and \(T_j\).

Assume explicitly that every earlier stage needed below has the two
lower-rung conclusions from (4.1): its vertex set is elementary under
its generating coloring (technical assertion (i)), and the
maximum-defective order assertion (v) has the form stated in
\noderef{ExtensionBaseSupportOrder}.  Then, for all (s,t\in L) with
(s\leq t),
\[
 \delta_t\notin\overline\varphi_{s-1}(T_s),
 \qquad \delta_t\ne\gamma_s.                                  \tag{1}
\]
Moreover, for every (t\in L),
\[
                  \delta_t\notin\overline\varphi_{j-1}(T_j).  \tag{2}
\]
The gamma and delta colors in Lean are tied to the exact optional fields
of their PE slots; neither conclusion is an additional hypothesis.
\end{lemma}
\begin{proof}
We first note a certified fact valid at every PE slot (r):
\[
                  \delta_r\notin
                  \overline\varphi_{r-1}(T_r).                 \tag{3}
\]
Indeed, eliminate the certificate at (r) using
\noderef{TashkinovSeries} and the equality
\(\Theta_r=\mathrm{PE}\).  In the PE constructor of
\noderef{TashkinovSeriesStep}, the selected coloring \(\pi_{r-1}\)
from \noderef{MaximumDefectiveChoice} is
((T_r,D_{r-1},\varphi_{r-1}))-stable, and its selected connecting edge
is a \(\delta_r\)-colored boundary edge.  If \(\delta_r\) were missing
at a vertex of (T_r) under \(\varphi_{r-1}\), the vertex clause of
\noderef{StableColoring} would make it missing there under
\(\pi_{r-1}\).  The old tree is
\(\varphi_{r-1}\)-closed by the PE certificate.  Stability preserves
both its vertexwise missing sets and every incident edge whose old color
lies in the controlling set, so the usual two applications of
\noderef{sc2} transfer this closedness to \(\pi_{r-1}\).  That contradicts
the displayed \(\delta_r\)-boundary edge.  This proves (3).

Fix (s,t\in L) with (s\leq t).  If (s=t), the first assertion of
(1) is (3).  The PE rule \noderef{ParallelExtensionColor} makes
\(\gamma_s\) missing at the supporting vertex under the selected
coloring, and stability makes it missing there under
\(\varphi_{s-1}\).  Hence
\(\gamma_s\in\overline\varphi_{s-1}(T_s)\), whereas (3) excludes
\(\delta_s\) from that set; therefore
\(\delta_s\ne\gamma_s\).

Suppose now that (s<t), and assume for a contradiction that
\(\delta_t\in\overline\varphi_{s-1}(T_s)\).  Choose
\(u\in V(T_s)\) at which \(\delta_t\) is missing.  Since (t\leq n)
and (s<t), this (T_s) is an earlier-rung stage, so the explicit
lower-rung assertion (i) says that its vertex set is elementary under
\(\varphi_{s-1}\).  By \noderef{ColoredVertexSetElementary}, no other
vertex of (T_s) misses \(\delta_t\).  Thus, for every
\(x\in V(T_s)\setminus\{u\}\), the definition
\noderef{MissingColors} supplies an incident edge (f_x) with
\(\varphi_{s-1}(f_x)=\delta_t\).

The PE certificate at (s) says that (T_s) is closed under
\(\varphi_{s-1}\).  Since \(\delta_t\) is missing at \(u\), closedness
\noderef{ColoredVertexSetClosed} excludes every
\(\delta_t\)-colored boundary edge.  Hence each exact identity (f_x)
is internal to (T_s), in the sense of \noderef{InternalEdges}.
Apply \noderef{SameColorSeriesPreserved} over the interval from (s) to
just before (t).  Its internal-edge clause gives
\[
                 \varphi_{t-1}(f_x)=\delta_t.                  \tag{4}
\]
The inclusions in \noderef{HkuMissingStoredMonotone}, iterated over the
same interval, put \(\delta_t\) in the missing-plus-stored controlling
set at (T_t).  Therefore the selected stability and then the second
stability fields of the PE certificate at (t) preserve (4).  Write
\(\pi'_{t-1}\) for that second stable coloring.  We have proved that
every (x\in V(T_s)\setminus\{u\}) is incident with an internal
\(\delta_t\)-edge under \(\pi'_{t-1}\).  Properness consequently
prevents any such (x) from also being a
\(\delta_t\)-boundary vertex.

By \noderef{ExtensionBaseSupportOrder}, whose explicit lower-rung
assertion (v) is among the present hypotheses,
\(v_t\preceq v_s\).  The certified ordered initial-segment relations
from \noderef{TashkinovSeriesStep} and
\noderef{TreeSequenceInitialSegment} imply that every vertex of (T_t)
preceding (v_t) lies in (T_s): indeed (v_s\in V(T_s)), and
\(v_t\preceq v_s\).

It remains to compare this with the maximum-defective choice at (t).
Its selected \(\delta_t\)-boundary set is defective.  The second stable
coloring controls \(\delta_t\), so its exact \(\delta_t\)-boundary set is
the same; this follows directly from the edge clause of
\noderef{StableColoring}, with symmetry supplied by \noderef{sc2}.
Thus it contains (v_t) and at least one other boundary vertex (y).
The maximality field of \noderef{MaximumDefectiveChoice}, applied to the
second stable coloring after forgetting the one extra controlled color,
gives (y\prec v_t).  Hence both (v_t) and (y) lie in (T_s).
If (u=v_t), then (y\in V(T_s)\setminus\{u\}) is a boundary vertex,
contrary to the preceding paragraph.  If (u\ne v_t), then (v_t)
itself gives that contradiction.  This proves
\(\delta_t\notin\overline\varphi_{s-1}(T_s)\).

The PE rule \noderef{ParallelExtensionColor} and selected stability make
\(\gamma_s\) a member of
\(\overline\varphi_{s-1}(T_s)\).  The exclusion just proved therefore
gives \(\delta_t\ne\gamma_s\), completing (1).

For (2), the finite set (L) is nonempty because it contains the
terminal PE index.  Choose its least member (s_0).  We have just shown
that \(\delta_t\notin\overline\varphi_{s_0-1}(T_{s_0})\) for every
(t\in L).  Between (j) and (s_0-1), no PE can be supported at a
vertex of (T_j): such a support would have first-appearance index at
most (j), hence would define an earlier member of (L).  Apply
\noderef{HkuUnsupportedVertexColors} separately at every vertex of
(T_j).  It identifies that vertex's missing-color set under
\(\varphi_{j-1}\) with its set under \(\varphi_{s_0-1}\).  The ordered
initial-segment fields identify (T_j) with the corresponding prefix of
(T_{s_0}).  Taking the union over its vertices gives
\[
 \overline\varphi_{j-1}(T_j)
   =\overline\varphi_{s_0-1}(T_j)
   \subseteq\overline\varphi_{s_0-1}(T_{s_0}).
\]
The already established exclusion at (s_0) now proves (2).
\end{proof}
